\chapter{Political Philosophy Is Useless for the Study of Servers}

\begin{center}
    Author: Moldypizzar0ll. Editor: retrooper. Single-blind peer reviewed. % clearly indicate non-peer reviewed articles!
\end{center}

\begin{flushright}
    IMSC: j(2)-ml-2

    DOI: \doilink{10.5281/zenodo.23702491}
\end{flushright}

% ============================================================================

\section*{Abstract}
Players of sandbox games in virtual worlds are often inclined to think that the same social phenomena that spring up in real life are applicable to player-to-player interactions in game; further, that these interactions are often governed by broader moral or ethical principles that reflect those of the real world. Thus, for many it makes sense to apply political philosophy to these games, since they appear to give us a window into how these systems spring up in the real world. Minecraft is perhaps the clearest example; because of its open-ended structure, players consider it to be an ideal environment for simulating and theorizing about political communities. However, I will argue that insofar as Minecraft presupposes ideal conditions which players must then constrain themselves, political philosophy can only engage in a descriptive, not normative, role in the study of communities on Minecraft servers.

\section{Introduction}
Minecraft servers are interesting in myriad ways, not least due to the fact that we tend to think that they provide insights into how people behave under peculiar circumstances. Insofar as we think this to be the case, we are also tempted to make prescriptions for how those people \emph{ought} to behave. To engage in a community is not just to engage in a descriptive project, in which the practices of the community are described and represented in detail, but also to engage in a normative one in which a certain level of institutional development and social coherence is needed to ensure that things are in order within the community. Surely there would be few who would say that these assumptions are unreasonable; and from them we may be tempted to derive the notion that our ways of understanding what prescriptions we are beholden to in the real world, especially those that pertain to the political sphere, may also apply to communities broadly; Minecraft servers being no exception. Hence, our tendency to apply political philosophy to virtual worlds.

I will argue here, however, that many of these intuitions are mistaken; that, in fact, most of the questions with which political philosophy deals are scarcely of importance for understanding what goes on in Minecraft servers (though the case may be made for other games, I will limit the scope of this argument to Minecraft). In essence, political philosophy aims to generate normative prescriptions that apply to political societies, i.e., prescriptions that apply both to authorities (if one believes the set of legitimate authorities to be nonempty), and to people following those authorities. These prescriptions are motivated by real-world constraints, while communities in Minecraft are not so constrained; the place of political philosophy in Minecraft is thus \emph{posterior}, rather than \emph{anterior}, to what people ought to do; for political philosophy usually only serves to present what is possible to the players in terms of the kind of situation they would like to find themselves in rather than as a legitimate, normatively-binding force that acts upon them. Whatever normatively binding forces actually exist, exist on the level of \emph{individuals} rather than on the level of a political community.

\section{The purview of political philosophy}
\subsection{Resource distribution}
In the first place, much of political philosophy is centered around how to distribute resources fairly. There are competing views, even prior to strictly political ones, as to how adequate distributions are to be achieved. Pure egalitarians argue that resources should be distributed equally across the population without exception. Luck egalitarians argue much the same and hold that differences in outcomes across individuals should be a function only of any given individual’s own choices. Prioritarians argue that benefits given to the least well-off should be weighed more than those given to those better-off.~\citep{Parfit(1997)} There are, of course, different ethical theories that provide justification for various political theories. The prioritarian, for example, may be motivated by the utilitarian consideration that one should aim to maximize the amount of good done for the greatest number of individuals, which means giving benefits to the least well-off. Others may not be moved by utilitarianism and instead view the consequences of an action simply as a reason, yet not a \emph{final} reason for doing something. In Anscombe’s example,~\citep{Anscombe(1967)} five people are stranded on one island and one person on another; you manage to spot both islands and decide to save the one person. The five people, according to Anscombe, do not have a claim to having been \emph{wronged}; for someone was indeed saved, and the rescuer did not decide \emph{not} to save anyone. She discusses another example, more explicitly related to the distribution of resources. One person suffers from a fatal illness, and a group of five suffer from the same illness. A single vial of an antidote is available; except the former person requires the entire vial to be cured, and the latter five each require one-fifth of the vial to be cured. No-one is wronged, Anscombe concludes, if one decides to give the one person the entire vial rather than saving the five. Is this a fair use of resources? Now, political considerations are, at bottom, ethical considerations. But there is disagreement with respect to how one moves from the ethical to the political. A plausible answer is that the ethical becomes the political when resources are involved on a sufficient scale---viz., when certain patterns of distribution may be considered more just than others. Whatever normative prescriptions follow, then, aim to achieve such distributions.

What we should note at this point is not strictly anything having to do with the individual views given here. Rather, we should note that the central aim of ethics is to remedy an existing state of affairs, call this \emph{A}, via some theory so as to achieve a fairer, more just state of affairs; call this \emph{B}. Different views have different verdicts on how to get from \emph{A} to \emph{B}, and what \emph{B} would actually look like. Some, such as Anscombe, do not believe \emph{B} to be wholly determined by the consequences of actions, and thus does not necessarily prescribe that \emph{B}, for a good ethical theory, always supervenes on consequences. Others, such as Parfit, would disagree.~\citep{Parfit(1984)} One of the plausible, though perhaps not exhaustive, ways in which a view can conceive this move from mere ethical prescriptions to talk of the political (taken in the philosophical sense) is in terms of the distribution of resources.

In Minecraft servers, however, things work differently. Assuming we examine a vanilla server in its infancy, its initial conditions are starkly different than any “initial” state of affairs in the real-world. Players are hidden behind a Rawlsian “veil of ignorance”, in which their physical real-world characteristics are largely stripped away and only their in-game capabilities and knowledge remain.\footnote{Of course, their mental capacities, memories, previous relationships, etc., still remain. What matters is that each player has the capacity to affect and derive resources from the world similarly, which is not so in the real world.}~\citep[p.~10]{Rawls(1971)} The distribution of resources, moreover, is nominally equal. The Rawlsian metaphor here breaks down, as we cannot benefit the least well-off having assumed the veil of ignorance if the distribution of resources past a certain point is scarcely of importance. Of course, there may exist a scarcity of \emph{labor};\footnote{I'm indebted to Snorka for this observation.} but there is little one cannot acquire with reasonable effort. Others have noticed this feature of the game, and many agree that the player is almost completely autonomous resource-wise. What are the possibilities, then, if one wants to make a server interesting? The traditional problems that drop out from the assumptions of real-world political philosophy do not really drop out in the case of Minecraft. Theoretically, each player can act autonomously resource-wise, and questions concerning how to distribute resources fairly do not arise. We are missing, however, the fact that this is precisely why political philosophy does not have any \emph{normative} pull in Minecraft, at least in the case of resource distribution. For the good server owner knows that an effective way to motivate interest in a server is to \emph{create} problems. The server owner may place borders on the world so as to limit the number of resources available; cap the number of valuable items in circulation; limit movement within a region to force cooperation; etc. It is only here that the traditional problems pop out. Indeed, if, for example, the world border is limited, one has to figure out the ideal distribution. But because such artificial constraints can only be generated by presupposing what the ideal distribution actually is (i.e., something like a server’s unaltered state where players are completely resource-autonomous), the solution is already baked into the design. While real-world political philosophy is indeterminate as to what the ideal state of affairs (\emph{B}) would look like from the perspective of a determinate state of affairs (\emph{A}), a Minecraft server has to \emph{create} indeterminacy by presupposing what the ideal conditions are and subsequently imposing constraints. Political philosophy, then, is only useful for people invested in the server insofar as the ones in charge of the server use it to \emph{design} the server, rather than provide normative prescriptions for those within it. At that point, however, it ceases being political philosophy.

Political philosophy gets its normativity precisely because the state of affairs towards which it aims is indeterminate and hazy. We \emph{ought} to do \emph{X} in part because the possibility of \emph{X} coming about is a \emph{possibility}, not an actuality. We cannot say that one \emph{ought} to do something given that the state of affairs with which they are charged with realizing has already come about.\footnote{N.B. This is distinct from normative prescriptions which say that one ought to \emph{repeat} certain actions so as to bring a certain state of affairs about continuously (e.g., continuously donating one tenth of your wealth to charity). The claim here pertains only to past states of affairs, which cannot be altered.} In server design, by contrast, the tools political philosophy has at its disposal can only be cashed out in terms of \emph{descriptions}. When the server owner imposes a world border, given what we know about what is valuable in the game, what the players ought to do is a foregone conclusion. We can only observe them going through the motions of playing within the world, knowing what the optimal decision procedures are for them already given the constraints. Now, the claim is not that the players are completely predictable; of course, players often do things they ought not to do (e.g., killing other players, griefing, hoarding, etc.), but that is not what is at stake here. The claim is that normative prescriptions cease being interesting as the \emph{raison d’etre} of a community when the structures that organize that community already presuppose what people ought to do in the first place. Political philosophy in this case becomes an \emph{explanation} for a state of affairs, rather than a guide for action so as to reach a different state of affairs.

\subsection{The state}
Take another salient topic in political philosophy: namely, the ideal position and configuration of the state. This is an issue with a distinguished history; its first true treatments may be found in Plato and Aristotle. The former thought the ideal form of political organization to consist of philosopher-rulers who, having been trained to seek truth and inquire into the nature of things (via the Forms), would be the best suited to manage a polity.~\citep{Plato(1968)} The latter (who was not terribly keen on the Forms) was concerned with the structure and size of the state as much as he was with what it ought to do. Should we have a monarch, in which only one individual rules, an aristocracy in which only a few rule, or a constitution\footnote{Aristotle used the term “democracy” to refer to the degenerate form of a constitution, in which the tyranny of the majority prevails. One may think of a “constitution” as something like a just, modern-day democracy. Additionally, for the three “just” versions of the possible configurations of the state, Aristotle posits an “unjust” mirror image of each of those forms; namely: \emph{tyranny} for monarchy, \emph{oligarchy} for aristocracy, and \emph{democracy} for the constitution.} in which the majority rule? \citet{Aristotle(2008)} thought a benevolent king to be the most ideal, though a constitution to be the most practical and just. Others were concerned with the function and limits of the state. Should the state wield unlimited power, or should it have checks and balances? Should we accept the latter, what kinds of checks should we impose? Hobbes was willing to accept the former;~\citep{Hobbes(1994)} others like Montesquieu and Nozick the latter. Sovereignty may be something that, by virtue of its nature, is indivisible---where any division in the power of the sovereign may lead us to suppose that whatever entities exist after the division are not sovereign at all. Given that an undivided, all-powerful sovereign is a preferable sort of political organization to an anarchic state of nature, we should opt for the former. This is the essence of Hobbes’s argument. Nozick is a good foil; for he argues that the sole permissible functions of the state are to enforce contracts and protect its subjects.~\citep{Nozick(1974)} Beyond this point, the state is morally unjustified in doing anything further. Others believe that the state should not exist at all, or at least that there are no legitimate authorities.

Again, the important thing to note here is that the normative prescriptions we put to states, whether it be with respect to its size, limits, or constitution, are a product of constraint. Political philosophers argue that we ought to move from one state of affairs to another, given the inadequacy of the preceding state of affairs. Minecraft, by contrast, can always \emph{be} adequate. We cannot prescribe how a state should look like in the game so much as one willingly and consciously participates in its construction. There is no \emph{normative} pull. If we want a particular sort of state to arise in the game, we must make the conditions suitable for its realization. We would not be concerned with what we ought to do in order to realize high-level goods, e.g., maximizing the satisfaction of the players, setting up the right kind of state, or distributing resources in the most efficient manner; for the best way to do that, presumably, is to let everyone run amok and do whatever they wish. Because of the game’s design as a sandbox, such ends are already built into the game. Minecraft is \emph{already} a maximizing tool, since, \emph{ceteris paribus}, it lets the player do whatever they want. Whatever political constraints there are appear to be artificial and scarcely normative. The place of political philosophy, then, if there is any, is in the realm of \emph{institutional design}; of designing certain situations so as to achieve certain objectives and outcomes. Its place is not in offering us any novelties about what we ought to do in certain situations, since the game, at least politically, already gives us solutions.

\section{Ethics, generally considered}
Now, in addition to this claim I also sought to claim something else; namely, that whatever normative prescriptions we \emph{can} cash out of Minecraft servers, occur on the level of individuals rather than political communities. More specifically, normative ethics as applied to individuals still holds up in the presence of changes in institutional design. We have observed that a political community in Minecraft must in large part be designed so as to achieve certain specifications and objectives; for such communities would not otherwise have arisen. Normative relations between individuals, however, are scarcely so designed. Although changes in an external situation can enforce certain changes on the entirety of a community, the relations between individuals do not necessarily change character as a result; normative prescriptions between them still have the same character. For example, if it is generally permitted on a server for players to kill each other, this does not necessarily make it morally permissible for individuals to kill each other in turn. I may still have an obligation to my friend to refrain from killing them as far as killing them would cause them harm in-game (and in certain exceptional circumstances, in the real world). But the institutional design of the server has no bearing whatsoever on whether my killing my friend is permissible or impermissible. What institutional design can do is make certain actions more or less \emph{acceptable}; for what matters in the case of institutional design is upholding the design itself, not necessarily maximizing a good. Hence breaking or following certain precepts may gain or lose one repute or respect. But this has no real normative import. If one desired to implement a monarchical system on a server, the question of whether the monarchy is, in fact, implemented can be pulled apart entirely from whether this arrangement benefits the players substantially. This is completely unlike the real world, wherein the maintenance of a state of affairs indeed has a direct bearing on the well-being of its participants.

To clarify: suppose that considerations of institutional design led us to desire the realization of a pattern, \emph{P}, within a server. We may also suppose that the purpose of a server is to complete some objective, \emph{O}, where \emph{O} takes into account in some measure the well-being of the participants, or some other ethically relevant aspect of the participants’ experience in the server. The claim is that \emph{O} is not necessarily dependent on the mere realization, or process of realizing \emph{P}. For the relation is actually inverted; depending on the objective of the server, the institutional design will reflect that objective. \emph{P}, in fact, depends on \emph{O}; and this dependence prefigures the relevant politico-philosophical precepts that account for what people ought to do once this pattern is established. Hence, any deviation from \emph{P} can be accounted for by the fact that it may just reflect \emph{O} more satisfactorily. Ethics \emph{in general}, as it applies to the individual player, \emph{can} figure in since that is precisely what determines the shape of \emph{O}. It is important not to conflate such prescriptions with politico-philosophical ones. It still makes perfect logical sense, for instance, to talk about the aggregate utility of a group of agents without invoking the state or authority in general, and consequentialist and nonconsequentialist ethics without patterns of resource distribution.

These conclusions are intuitively absurd. However, I am not claiming that political philosophy is useless in Minecraft in general. Rather, I am claiming that it is not useful when we consider the study of human behavior in Minecraft servers specifically. Of course, were Minecraft like the real world to a high degree of accuracy, it would certainly be useful. But presumably the study of these pockets of human interaction is not the study of what individuals \emph{ought} to do, but rather of \emph{what} they do. What they ought to do can only be cashed out at the level of the individual, as far as they are charged with the presence of other players with whom they interact. But when it comes to the general design of the server in which they find themselves, any large-scale normative prescriptions that could hope to approximate anything that real-world political philosophy would offer vanish---and for the simple reason that one can theoretically realize whatever state of affairs they want within the game.

\section{Concluding thoughts}
What we find interesting about applying political philosophy to Minecraft is not, as our intuitions may tell us enticingly, the fact that we find moral prescriptions for what people should do under certain circumstances. Rather, what we find interesting is isolating certain facets of human behavior that we know to be present \emph{already} and seeing them play out in real time without real-world consequences. Political philosophy cannot provide anything for the game for precisely this reason. When designing servers, or places where people coexist in general, one must inevitably come up against a form of the old Hegelian problem that a movement of a state of affairs is not necessary but \emph{imposes} its necessity. In Minecraft servers, the ethical never becomes the political; for the political is imposed, not made.

\section{Declarations}
No conflict of interest and no AI use have been declared.

    {
        \raggedright
        \begin{thebibliography}{99}
            \bibitem[Anscombe(1967)]{Anscombe(1967)}
            Anscombe, G. E. M. “Who Is Wronged? Philippa Foot and ‘Double Effect.’” The Oxford Review, vol. 5, Oct. 1967, pp. 16–17.
            \bibitem[Aristotle(2008)]{Aristotle(2008)}
            Aristotle. The Basic Works of Aristotle. Edited by Richard McKeon and Richard McKeon, Random House, 2008.
            \bibitem[Hobbes(1994)]{Hobbes(1994)}
            Hobbes, Thomas. Leviathan. 1668. Edited by Edwin Curley, Hackett, 1994.
            \bibitem[Nozick(1974)]{Nozick(1974)}
            Nozick, Robert. Anarchy, State, and Utopia. Basic Books, 1974.
            \bibitem[Parfit(1997)]{Parfit(1997)}
            Parfit, Derek. “Equality and Priority.” Ratio, vol. 10, no. 3, Dec. 1997, pp. 202–21, https://doi.org/10.1111/1467-9329.00041.
            \bibitem[Parfit(1984)]{Parfit(1984)}
            Parfit, Derek. Reasons and Persons. Oxford University Press, 1984.
            \bibitem[Plato(1968)]{Plato(1968)}
            Plato. The Republic of Plato. 1968. Translated by Allan Bloom, Basic Books, 1968.
            \bibitem[Rawls(1971)]{Rawls(1971)}
            Rawls, John. A Theory of Justice. Harvard University Press, 1971.

        \end{thebibliography}
    }